\documentclass[10pt,a4paper]{amsart}
\usepackage[margin=19mm]{geometry}
\usepackage{amsmath,amssymb,mathtools}
\usepackage[scr=rsfs,cal=euler]{mathalfa}
\usepackage{xfrac}
\usepackage[unicode,hidelinks,bookmarks=false]{hyperref}
\newcommand{\ie}{i.e.\ }
\newcommand{\cf}{cf.\ }
\newcommand{\resp}{respectively\ }
\newcommand{\Subsection}{\subsection}
\providecommand{\nobreakdash}{\nobreak\hbox{-}}
\setlength{\emergencystretch}{2em}
\multlinegap0pt
\usepackage{bm}
\usepackage{bbm}
\usepackage{dsfont}
\usepackage{xspace}
\usepackage{cancel}
\usepackage{mathtools}
\usepackage{amsthm}
\makeatletter
\@ifundefined{theorem}{\newtheorem{theorem}{Theorem}}{}
\makeatother
\newcommand{\C}{\mathbb{C}}
\title[Lifts of Operator Systems]{Lifts of Operator Systems}
\author{Markus Dannem\"uller, Tim Netzer}
\date{Source: arXiv:2508.15348v1 (CC BY 4.0)}
\begin{document}
\maketitle

\noindent\emph{Source and licence.} Lifts of Operator Systems. Version arXiv:2508.15348v1 (CC BY 4.0), distributed under the Creative Commons Attribution 4.0 International licence, as recorded in the arXiv licence metadata for this exact version and shown on the abstract page https://arxiv.org/abs/2508.15348v1. Full rights notice: licenses/arxiv-2508.15348-CC-BY-4.0.txt.\par\smallskip

\noindent\textit{Editorial scope.} Counted unit: the Theorem whose source label is \texttt{thm:NC\_lift\_iff\_fact} in the article (source0.tex lines 151--153, source label \texttt{thm:NC\_lift\_iff\_fact}), delivered together with the article's own complete original proof of it (source0.tex lines 154--171, one \texttt{proof} environment of 18 lines). No mathematics is written, repaired or re-derived: statement and proof are the article's own text line for line. Required context retained uncounted: none. Every definition, display and label of those lines is retained unchanged. Omitted from this package and not redistributed: 1--150, 172--288. Nothing inside a retained excerpt is omitted or altered: every retained body is delivered exactly as the article's lines are written, so no omission receipt of any kind accompanies this package. Citations: the retained excerpts carry no citation command at all, so no bibliography entry of the article is transcribed and no reference is redistributed. Reference adaptation: none was needed and none was made; every reference inside the delivered unit resolves inside it, so no unresolved reference or citation is produced. Printed slips kept literal: no printed word, symbol, hypothesis or number of the counted statement or of its proof was altered. Layout and class: the article's own class is \texttt{article} with packages including amsmath, amssymb, amsthm, graphicx, mathtools, authblk, hyperref, cleveref, tikz-cd; the wrapper is the lane's pinned \texttt{amsart} editorial wrapper at 10pt on A4, which restates the theorem-like environments the retained text needs and adds the article's own macro definitions that the retained text needs (\texttt{\textbackslash C}), with extra wrapper packages (bm, bbm, dsfont, xspace, cancel, mathtools). The delivered numbering is the unit's own. Assets: no retained span contains a figure environment, a graphics-inclusion command, a PDF-page inclusion, a table or a TikZ picture; no image file is copied into this package, the package declares no asset dependency and no image file is redistributed with it. Licence: version arXiv:2508.15348v1 of this article is distributed under the Creative Commons Attribution 4.0 International licence, as recorded in the arXiv licence metadata for this exact version and shown on the abstract page https://arxiv.org/abs/2508.15348v1. A full rights notice is delivered at licenses/arxiv-2508.15348-CC-BY-4.0.txt. Characters: every retained excerpt is pure ASCII, so the delivered engine, XeLaTeX with its default Latin Modern text font, reports no missing character.
\begin{theorem} \label{thm:NC_lift_iff_fact}
	Let \(S\) and \(T\) be operator systems on \(X\) and \(Y\), respectively. If \(S\) admits a proper \(T\)-lift, then its slack operator  admits a \(T\)-factorization. Conversely, if the slack operator of $S$ admits a \(T\)-factorization, then \(S\) admits a  \(T\)-lift.
\end{theorem}

\begin{proof}
	Let us assume first that \([\cdot,\cdot]_S\) admits a \(T\)-factorization, given by the family of maps \((\alpha, \beta)\). We define the vector space \[Z \coloneqq \lbrace (x, y) \in X \oplus Y \mid \forall t, \phi \in S^\lor_t\colon \phi(x) = \beta_t(\phi)(y) \rbrace,\] which comes equipped with the canonical projections \[\pi_X\colon Z \to X, \qquad \pi_Y\colon Z \to Y.\] Our goal is to show that \((\pi_X, \pi_Y)\) is a \(T\)-lift for \(S\). Hence, we need to show that $\pi_Y$ is injective and \(\pi_X[\pi_Y^{-1}[T]] = S\) holds. 

   For injectivity of $\pi_Y$, assume \((x, y) \in \operatorname{ker}(\pi_Y)\), so clearly \(y = 0\). But then \((x, 0)\) is such that \(\forall t, \phi \in S^\lor_t\) we have \(\phi(x) = \beta_t(\phi)(0) = 0\), since \(\beta_t(\phi)\) is a linear map. Already from the case $t=1$ we get  $x\in S_1\cap -S_1=\{0\}$, proving injectivity. 
	
    Now let \(c \in \pi_Y^{-1}[T]_s\), \(\phi \in S^\lor_s\), and consider \(\phi[\pi_X[c]]\in{\rm Her}_{s^2}(\C).\) By construction, the maps \(\phi \circ \pi_X\) and \(\beta_s(\phi) \circ \pi_Y\) coincide on \(Z\). But then they  also agree when applying them to elements of \(\pi_Y^{-1}[T]_s\), so we get \[\phi[\pi_X[c]] = \beta_s(\phi)[\pi_Y[c]].\] Since \(\pi_Y[c] \in T_s\) and \(\beta_s(\phi) \in T^\lor_s\), the right-hand side is positive semi-definite, and this in particular implies  $\phi(\pi_X[c])\geqslant 0.$ Since this holds for every \(\phi \in S^\lor_s\), we get $\pi_X(c)\in S_s$, by classical biduality. This shows that \(\pi_X[\pi_Y^{-1}[T]] \subseteq S\).
	
	For the reverse inclusion, let \(a \in S_s\). Since \(\alpha_s(a) \in T_s\), we clearly have $\pi_Y^{-1}[\alpha_s(a)] \subseteq \pi_Y^{-1}[T]_s.$ We now define $$\tilde{a} \coloneqq (a, \alpha_s(a))\in (X\otimes{\rm Her}_s(\C))\oplus(Y\otimes{\rm Her}_s(\C))=(X\oplus Y)\otimes{\rm Her}_s(\C).$$ Once we have shown that $\tilde a\in Z\otimes{\rm Her}_s(\C)$ holds, we are done, since  \(\tilde{a} \in \pi_Y^{-1}[\alpha_s(a)]\subseteq \pi_Y^{-1}[T]_s,\) and $\pi_X[\tilde a]=a$ holds. To prove that $\tilde a$ lives over $Z$, let us fix a basis \((M_i)_i\) of \(\operatorname{Her}_s(\C)\) and write \[a = \sum_i x_i \otimes M_i, \qquad \alpha_s(a) = \sum_i y_i \otimes M_i.\]  From the definition of a factorization we obtain \begin{align*}\sum_i\beta_t(\phi)(y_i)\otimes M_i&=\beta_t(\phi)[\alpha_s(a)]=\phi[a]=\sum_i\phi(x_i)\otimes M_i.\end{align*} Since the $M_i$ form a basis, this implies $\beta_t(\phi)(y_i)=\phi(x_i)$ for all $i$, i.e.\ $(x_i,y_i)\in Z.$ Now clearly $\tilde a=\sum_i (x_i,y_i)\otimes M_i\in Z\otimes{\rm Her}_s(\C)$, as desired.
	
	Taking both inclusions together, we conclude that \(S = \pi_X[\pi_Y^{-1}[T]]\), hence \((\pi_X, \pi_Y)\) forms a \(T\)-lift for \(S\).
	
	For the converse statement, let us assume that \(S\) admits a proper \(T\)-lift, given by  linear maps \(\pi\colon Z \to X\) and \(\gamma\colon Z \to Y\) for some vectorspace \(Z\), such that $\gamma$ hits the interior of $T_1$. 
	
	We note that for any \(a \in S_s\), the fibre \(\pi^{-1}[a]\) has non-empty intersection with \(\gamma^{-1}[T]\). So we can just pick an arbitrary  element \(\tilde{a} \in \pi^{-1}[a] \cap \gamma^{-1}[T]\) and define \(\alpha_s(a) \coloneqq \gamma[\tilde{a}]\in T_s.\)
	
	Next we need to construct \(\beta_t(\phi)\in T^\vee_t\) for any given \(\phi \in S^\lor_t\). Understanding again \(\phi\) as an element of \({\rm CP}(S, P^t)\), we pull this map back along \(\pi\) and define \[\tilde{\phi} \coloneqq \phi \circ \pi.\] For any \(c \in \gamma^{-1}[T]\) we have \[\tilde{\phi}[c] = \phi[\pi[c]],\] and since \(\pi[c] \in S\) and \(\phi \in S^\lor\), this expression is positive semidefinite. This shows that \[\tilde{\phi} \in {\rm CP}(\gamma^{-1}[T], P^t).\]
	We now  extend \(\tilde{\phi}\) to a completely positive map on \(T\), which we then take as \(\beta_t(\phi)\). This  makes sense since \(\gamma\) is injective, and we can thus view it as the inclusion of \(\gamma^{-1}[T]\) into \(T\), and  Arveson's Extension Theorem guarantees the existence of an extension (which uses properness of $\gamma$).   We now have  \[[\alpha_s(a),\beta_t(\phi)]_T=\beta_t(\phi)[\alpha_s(a)] = \tilde{\phi}[\tilde{a}] = \phi[\pi[\tilde{a}]] = \phi[a]=[a,\phi]_S,\] which shows that \((\alpha, \beta)\) is indeed a \(T\)-factorization of the slack operator of \(S\). 
\end{proof}

\end{document}
